% Generated draft: requires mathematical and source review.
\documentclass[11pt]{article}
\usepackage[margin=25mm]{geometry}
\usepackage{fontspec}
\setmainfont{Latin Modern Roman}
\usepackage{amsmath,amssymb,mathtools,unicode-math}
\setmathfont{Latin Modern Math}
\usepackage{enumitem,needspace,xurl,hyperref}
\setlength{\parindent}{0pt}
\setlength{\parskip}{5pt}
\setlength{\emergencystretch}{3em}
\newcommand{\N}{\mathbb{N}}
\newcommand{\Z}{\mathbb{Z}}
\newcommand{\Q}{\mathbb{Q}}
\newcommand{\R}{\mathbb{R}}
\newcommand{\C}{\mathbb{C}}
\newcommand{\F}{\mathbb{F}}
\newcommand{\A}{\mathbb{A}}
\newcommand{\PP}{\mathbb{P}}
\newcommand{\E}{\mathbb{E}}
\newcommand{\eps}{\varepsilon}
\newcommand{\dd}{\,\mathrm{d}}
\begin{document}
\section*{UnsolvedMath ID: 30006331}
\textbf{Source problem code:} catalogue entry
\subsection{Definitions and mathematical statement}
The supplied record asks about symmetry methods for geometric variational problems involving polygons, but does not specify the energy, the original admissible polygon class, the mass constraint, or the anisotropy defining the Wulff shape. Consequently, it does not determine a self-contained conjecture with fixed hypotheses.

To describe the missing formulation, let \(\mathcal A_{m,n}\) provisionally denote an admissible class of polygons of mass \(m>0\), with a side-count parameter \(n\), and let \(E\) denote the energy to be studied. These symbols are placeholders: the record does not define mass, the admissible polygons, whether the side count is exact or bounded, or the domain and formula of \(E\). Once these data are fixed, a minimizer \(P_*\in\mathcal A_{m,n}\) would satisfy
\[
E(P_*)\leq E(P)\qquad\text{for every }P\in\mathcal A_{m,n}.
\]
If criticality is defined using differentiable admissible variations, a critical polygon \(P\) would satisfy
\[
\left.\frac{\mathrm d}{\mathrm dt}E(P_t)\right|_{t=0}=0
\]
for every allowed differentiable path \((P_t)_{|t|<\delta}\subseteq\mathcal A_{m,n}\) with \(P_0=P\), whenever the derivative exists. The source does not specify the allowed variations or establish that this differentiable notion is the intended one.

\textbf{Source question, with unresolved qualifications.} Determine how the symmetry-based minimization techniques from the unspecified original polygon problem extend to critical points and to polygons with a larger number of sides. Determine an appropriate meaning of a ``most symmetric object'' in the enlarged class. Then determine whether the associated Wulff shape minimizes the intended energy for sufficiently small masses.

If a mass-normalized Wulff competitor \(W_{m,n}\in\mathcal A_{m,n}\) is defined by the missing data, the final question has the schematic form
\[
\exists\,m_0>0\quad\forall\,m\in(0,m_0)\quad
\forall\,P\in\mathcal A_{m,n},\qquad
E(W_{m,n})\leq E(P).
\]
This display describes a possible logical form, not a fully specified conjecture: the range of \(n\), the dependence of \(m_0\), the relation between symmetry and anisotropy, and even admissibility of the Wulff competitor remain unresolved. The source does not specify uniqueness or the behavior at \(m=m_0\).
\subsection{Short English statement}
The problem asks whether symmetry methods for polygonal minimization also describe critical points and polygons with more sides, and whether a suitable Wulff shape minimizes the energy at sufficiently small mass. The supplied excerpt omits the energy, admissible class, mass definition, and anisotropy, so a precise conjecture requires the missing context. The supplied status is \emph{partially solved}; it is a dated curation label, not an independently verified conclusion.
\subsection{Sources}
\begin{itemize}
\item Supplied local UnsolvedMath registry record, permanent ID \(30006331\), title \emph{Critical Points of Symmetric Geometric Energies}, source problem code \texttt{catalogue entry}, in the Oberwolfach Reports Open Problems collection. The record gives proposed year 2025 and does not identify a proposer.
\item \emph{Singularities in Discrete Systems} (2025), Oberwolfach Reports, pp.\ 1163--1214, DOI \texttt{10.4171/owr/2025/23}. Supplied source report and literature reference: \url{https://doi.org/10.4171/owr/2025/23}. The record describes these reports as documenting mathematical workshops and their open-problem discussions; it provides no page locator for this particular question.
\item \emph{Minimality of Polytopes}. Supplied literature reference: \url{https://doi.org/10.1016/j.na.2020.112223}.
\item \emph{Riesz-type inequalities for polygons}. Supplied literature reference: \url{https://arxiv.org/abs/2103.06657}.
\item Literature assessment dated \texttt{2026-08-22}: the supplied curation labels the problem partially solved and states that the primary work establishes restricted symmetric polygon and small-mass cases, without a general classification or minimizer theorem covering critical points, arbitrary numbers of sides, and the proposed Wulff-shape conclusion. This is supplied dated curation derived from the OWR material, not an independently verified literature assessment. Its time-sensitive status and cited claims require independent verification.
\item The registry labels the statement \texttt{corrected\_verified} and reports that the official report was checked, the question isolated from surrounding exposition, and its notation repaired. These are supplied provenance claims; this draft does not independently verify them. The original-statement field describes its passage as retained verbatim, but the surrounding mathematical setup is absent from the supplied record.
\end{itemize}
\end{document}
